\begin{definition}[Extractors and dispersers]
Let $\Gamma$ and $\Omega$ be some finite domains.
Let ${\cal C}$ be a class of random variables taking values in  $\Gamma$.
We say a random variable $P$ taking values in $\Omega$ is
\emph{$\eps$-close to uniform} if for every  $A\subseteq \Omega$,
%
%
\[\left|\Pr (P \in A) - \frac{|A|}{|\Omega|}\right| \leq \eps\,.\]
\begin{itemize}
%
%
%
%
\item Fix any $0\leq \eps <1$. A function $E:\Gamma \to\Omega$
is an \emph{$\eps$-extractor for ${\cal C}$}
if for every $X \in {\cal C}$,
the random variable %
$E(X)$ is $\eps$-close to %
uniform.
%
%
\item A function $D:\Gamma \to\Omega$ is
%
a \emph{disperser for ${\cal C}$} if for every $X \in {\cal C}$,
the random variable %
$D(X)$ takes more than %
one value in $\Omega$ with nonzero probability.
\end{itemize}
\end{definition}
\subsection{Polynomial sources}

%
%
%
In this paper we construct extractors for \emph{polynomial sources},
which are distributions that are sampled by applying low-degree
%
polynomials to uniform inputs as defined next.
%
%
%
Throughout this paper, if $\Omega$ is a finite set, we let
$U_\Omega$ denote the uniform distribution over $\Omega$.
%
%
By the \emph{individual degree} of a multivariate polynomial $f$ we
mean the smallest $d$ such that $f$ has degree $\le d$
in each variable.


\begin{definition}[Polynomial sources]\label{dfn-polysource}
Fix integers $n,k,d$ with $k\leq n$ and a field $\F$.
We define $\M{n}{k}{d}$ to be the set of mappings $f:\F^r\to\F^n$,
where $r$ is an integer counting the number of inputs to the source and
\[f(Z_1,\ldots,Z_r)=(f_1(Z_1,\ldots,Z_r),\ldots,f_n(Z_1,\ldots,Z_r))\] such that
\begin{itemize}
\item for every $i\in [n]$, $f_i$ is a polynomial in $\F[Z_1,\ldots,Z_r]$
of individual degree at most $d$.
\item The range
%
%
%
%
%
%
%
of $f$ is of size at least $q^k$. Formally,
\[|\{f(z_1,\ldots,z_r)\mid (z_1,\ldots,z_r)\in \F^r\}| \geq q^k\,.\]
\end{itemize}
%
A \emph{$(n,k,d)$-polynomial source} is a random variable of the form
$f(U_{\F^r})$ for some and $f\in \M{n}{k}{d}$ with $r$ inputs.
(When the parameters $n,k,d$ are clear from
%
%
the context, we shall omit them, and simply use the term
``polynomial source.'')

\end{definition}

%
%
\begin{definition}[Polynomial-source extractors]\label{dfn-polyext}
%
Let $\Omega$ be some finite set.
A function $E:\F^n\to \Omega$  is a \emph{\polyext{k}{d}{D}{\eps}} if for
every $f\in \M{n}{k}{d}$ of total degree at most $D$ and $r$ inputs,
$E(f(U_{\F^r}))$ is $\eps$-close to uniform
(where $U_{\F^r}$ denotes the uniform distribution
%
over $\F^r$).
%
%
%
%
%
\end{definition}
\begin{remark}
A few words are in order regarding \expref{Definition}{dfn-polysource}.
\begin{itemize}
\item
%
%
%
 The number of inputs used by our source, denoted by $r$ in
 \expref{Definition}{dfn-polysource}, does not affect the
%
%
parameters of our extractors, and hence
%
we omit this parameter from the definition of
polynomial sources and extractors.

\item In the context of extractors what might have seemed
more natural is to require the
%
%
%
%
random variable
$f(U_{\F^r})$ to have
\emph{min-entropy}\footnote{The min-entropy of a
%
%
%
random variable $X$ is
the largest %
real number $\ell$
such that for every
%
value $x$
we have %
%
%
%
$\Pr(X=x)\leq 2^{-{\ell}}$.
This is the standard measure of randomness in the context of extractors,
%
originating from Chor and Goldreich~\cite{CG88}.}
at least $k\cdot \log q$.
%
%
%
Our requirement on the size of the range of $f$ is seemingly weaker,
and suffices for our construction to work.
%
%
%
%
(In particular, our result implies that, for some settings of field order and degree, when $f$ has large range
the random variable $f(U_{\F^r})$ is statistically close to a random variable that has at least a certain min-entropy.)

\item Individual degree plays a larger role than total degree in our results.
%
  In fact, the first stage of our construction---constructing a
%
%
  non-constant polynomial over $\F$---requires a field of order
%
  depending \emph{only} on individual degree. This is why it is
  more convenient to limit individual degree and not total degree
  in the definition of $\M{n}{k}{d}$.

\end{itemize}
\end{remark}

\begin{theorem}[Main--Extractor]\label{thm-ext_main_intro}
Fix a field $\F$ of characteristic $p$, integers $d,D,4\leq k\leq n$ where $n\geq 25$,
and a positive integer $m<1/2\cdot \log_p q$.
Let $\alpha=3 D \cdot (p\cdot d)^{3 n/k} $. Assume that   $q\geq 2\cdot  \alpha^2$.
There is an explicit \polyext{k}{d}{D}{\eps}  $E:\F^n\to \mathbb{F}_p^m$ with error
$\eps =  p^{m/2}\cdot \alpha \cdot q^{-1/2}$.
\end{theorem}

\section{Overview of the proof}

Our goal is to describe an explicit function $E:\F^n \to \{0,1\}^m$
such that for any $(n,k,d)$-polynomial source $X$ we have that
$E(X)$ is $\eps$-close to the
%
%
uniform distribution over $\{0,1\}^m$.  We do this in two steps.

First we construct a function $E_0$, called a \emph{non-Boolean disperser}, that is guaranteed to be non-constant on $X$, \ie,
such that the 
%
%
random variable $Y=E_0(X)$ takes more than one value.
This part is done in \expref{Section}{sec:main}.
Then we apply a second function $E_1$ to the output of $E_0$ and prove, using the fact that $E_0$ is a low-degree function in our case, 
that the distribution of %
$E_1(Y)=E_1(E_0(X))$ is $\eps$-close to uniform.
 This ``disperser--to--extractor'' part is described in Sections~\ref{sec-weil} and~\ref{sec:extractor}.
 We now informally describe the two functions assuming for simplicity
that %
the field $\F$ is of characteristic $2$ and that $n$ is prime.
Before starting let us recall the notion of a Frobenius automorphism.
If $\K$ is a finite field of characteristic $2$ then the mapping
\[\sigma_i:\K\to \K, \qquad \sigma_i(z)=z^{2^i}\] is a
\emph{Frobenius automorphism of $\K$ over $\Ft$}.
%
%
%

%
%
%
%
%
%
%
%

%
%
%
%
The three elementary properties of this mapping that we use below are
%
%
%
%
%
%
%
%
%
%
\begin{itemize}
\item[(i)] its \emph{$\Ft$-linearity:} $\sigma_i(a+b)=\sigma_i(a)+\sigma_i(b)$,
\item[(ii)] its \emph{distinctness}: if $\K$ is an extension of $\Ft$ of
   degree at least $t$ and $0\leq i<j\leq t-1$ then $\sigma_i$ and
   $\sigma_j$ are different, and
\item[(iii)] its \emph{dimension-preservation:} If $\K\supset\F\supset\Ft$
 then $A\subset \K$ and $\sigma_i(A)\triangleq \{\sigma_i(a) \ | \ a\in A\}$
 span spaces of equal dimension over $\F$
 (see \expref{Claim}{claim-frobenius_has_dim}).
\end{itemize}

\paragraph{A different view of low-degree sources} The first part of
our analysis uses a somewhat nonstandard view of low-degree sources
that we need to highlight. The random variable $X$ ranges over $\F^n$
and is the output of $n$ degree-$d$ polynomials over $\F$. Let
\[
\monom
\]
denote the set of monomials over $\F$
of individual degree at most $d$ where $d<q$. (We use $Z$ variables to denote inputs of the polynomial source and $X$ variables for its output.) Suppose
%
the $i$-th
coordinate of $X$ is
\[X_i=P^{(i)}(Z_1, \ldots, Z_r)=\smashoperator[r]{\sum_{M\in \monom}} a^{(i)}_M \cdot M(Z_1,\ldots,Z_r) \]
where $a^{(i)}_M\in \F$ and $Z_1,\ldots, Z_r$ are independent random variables distributed uniformly over $\F$. Applying an $\F$-linear bijection $\phi:\F^n\to\Fqn$, let $a_M=\phi(a^{(1)}_M,\ldots,a^{(n)}_M)$ denote the sequence of coefficients of the monomials $M$, viewed now as a single element in $\Fqn$. Our nonstandard view is that our source is
\begin{equation}
  \label{eq:nonstandard} X=P(Z_1,\ldots, Z_r)= \smashoperator[r]{\sum_{M\in \monom}} a_M\cdot M(Z_1,\ldots,Z_r)
\end{equation}
where the coefficients $a_M$ and the random variable $X$ come from the ``large'' field $\Fqn$ but the random variables $Z_1,\ldots,Z_r$ still range over the ``small'' field $\F$. This large-field-small-field view will be important in what comes next. In particular, we shall use the following claim which reduces the problem of constructing a non-Boolean disperser to that of constructing a polynomial whose coefficients span $\Fn$ over $\F$.

\begin{claim}[Full-span polynomials are non-constant coordinate-wise]\label{clm:nonconstant}
Suppose $P$ has individual degree smaller than $q$.
If the set of coefficients $A=\condset{a_M}{\deg(M)>0}$ appearing in (\ref{eq:nonstandard}) spans $\Fn$ over $\F$
then $X_i=P^{(i)}(Z_1,\ldots,Z_r)$ is a non-constant function on $\F^r$ for every $i\in\set{1,\ldots,n}$.
\end{claim}

\begin{proof}
  By way of contradiction. If $P^{(i)}$ is constant on $\F^r$
and has individual degrees smaller than $q$, then as a formal polynomial
it is constant. This implies that all elements of $A$, as vectors in $\F^n$, are equal to zero in the
%
$i$-th coordinate.
Thus, $A$ spans a strict subspace of $\Fn$ in contradiction to the assumption of the claim.
\end{proof}

\paragraph{Non-Boolean disperser} We start with the simplest nontrivial case to which our techniques apply and construct a non-Boolean disperser for homogeneous multilinear quadratic sources with min-entropy rate greater than half over the finite field with $4$ elements (this is a special case of Corollary~1.8).
Using
%
%
$\binom{[r]}{2}$
to denote the set $\condset{(i,j)}{1\leq i< j\leq r}$ and writing $X$ as in (\ref{eq:nonstandard}) we get
\begin{equation}
  \label{eq:quadratic-multilinear}
  X=\smashoperator[r]{\sum_{(i,j) \in\binom{r}{2}}} a_{ij} Z_i Z_j, \qquad a_{ij} \in \Ffn
\end{equation}
where $Z_1,\ldots,Z_r$ are uniformly and independently distributed over 
$\Ff$ and $X$ 
%
%
takes more than $4^{n/2}$ distinct values.
Let \begin{equation}
  \label{eq:A}
  A=\left\{ a_{ij} \;\middle|\; (i,j) \in \binom{[r]}{2}\right\}
\end{equation}
denote the set of coefficients appearing in (\ref{eq:quadratic-multilinear}). In light of \expref{Claim}{clm:nonconstant} it suffices to construct $E_0$ such that $E_0(X)$,
 when written as a polynomial over $Z_1,\ldots, Z_r$, has a set of coefficients that spans $\Ffn$ over $\Ff$. (Then we ``project'' this polynomial onto, say, the first coordinate and get a non-constant function mapping into $\Ff$, \ie, a non-Boolean disperser.)

To do this we take the approach of DeVos and Gabizon~\cite{DG10} which uses the theorem of Hou, Leung and Xiang~\cite{HLX02}.
 Assuming $n$ is prime, this theorem implies that if $A,B\subset \mathbb{F}_{q^n}$ are sets spanning spaces of respective dimensions $d_1,d_2$ over $\F$,
 then the set of products \[A\cdot B\triangleq \condset{a\cdot b}{a\in A, b\in B}\] spans a subspace of $\mathbb{F}_{q^n}$ over $\F$ of dimension at least $\min\{n, d_1+d_2-1 \}$. Returning to our case and taking $A$ as in (\ref{eq:A}), our first observation is that $\dim(\spn(A))>n/2$ because $X$ is contained in $\spn(A)$. So the theorem of~\cite{HLX02} mentioned above implies that $\spn(A\cdot A)=\Ffn$.
Consider what would happen if we could sample \emph{twice} from $X$ independently and take the product of the two samples in $\Ffn$. Using $X',Z'_1,\ldots,Z'_r$ to express the second sample we write this product as
\begin{align*}
  X\cdot X' &=  \left(\smashoperator[r]{\sum_{(i,j) \in\binom{r}{2}}} a_{ij} Z_i Z_j \right)\cdot \left(\smashoperator[r]{\sum_{(i',j') \in\binom{r}{2}}} a_{i'j'} Z'_i Z'_j \right).
\end{align*}
Opening the right-hand-side as a polynomial in $Z_1,\ldots,Z_r,Z'_1,\ldots,Z'_r$ we see that its set of coefficients is $A\cdot A$ which spans $\Ffn$ over $\Ff$, as desired.\footnote{The same argument would work as well over the two-element field $\mathbb{F}_2$. The extension field is needed to deal with the case of a single source as explained next.}


Unfortunately we only have access to a \emph{single } sample of $X$ and have to make use of it. We use the fact that $\Ff$ is a degree $2$ extension of a smaller field ($\mathbb{F}_2$) and hence has two distinct Frobenius automorphisms. And here comes our second observation: Taking the product of $2$ distinct Frobenius automorphisms of a \emph{single} sample of $X$ has a similar effect to that of taking two independent samples of $X$! Indeed, take the product of $\sigma_0(X)$ and $\sigma_1(X)$ and, using the linearity of Frobenius mapping, expand as
\begin{align*}
X\cdot X^2 &= \left(\smashoperator[r]{\sum_{(i,j) \in\binom{r}{2}}} a_{ij} Z_i Z_j \right)\cdot\left(\smashoperator[r]{\sum_{(i', j')\in\binom{r}{2}}} a_{ij}^2 Z_i^2 Z_j^2 \right) \\
&= \smashoperator[r]{\sum_{(i,j), (i', j')\in\binom{r}{2}}} a_{ij} a_{i'j'}^2 Z_i Z_j Z_{i'}^2 Z_{j'}^2\,.
\end{align*}

The main point is that every element in the set of products of $A$ and $A^2\triangleq\condset{a^2}{a\in A}$ appears as the coefficient of a monomial in the polynomial above and these monomials are distinct over $\Ff$. And the dimension-preservation of $\sigma_1$ implies that $\dim(\spn(A^2))=\dim(\spn(A))>n/2$. Consequently, the theorem of~\cite{HLX02} implies that $A\cdot A^2$ spans $\Ffn$ over $\Ff$, so by \expref{Claim}{clm:nonconstant} the function $E_0(X)$, which outputs the first coordinate of $X\cdot X^2$, is non-constant for $X$ and this completes the sketch of our non-Boolean disperser for the special case of homogenous, quadratic, multilinear polynomials over $\Ff$.
%
%
%

To extend this argument to general polynomial sources of individual degree $\leq d$ we carefully select a set of $t$ distinct Frobenius automorphisms $\sigma_{i_0},\ldots, \sigma_{i_{t-1}}$ (assuming $\F$ is an extension-field of degree at least $t$) such that the
 mapping $f:(\monom)^t\to \F[Z_1,\ldots,Z_r]$ given by \[f(M_0,\ldots,M_{t-1})=\prod_{j=0}^{t-1} \sigma_{i_j}(M_j)\quad \bmod (Z_1^q-Z_1,\ldots, Z_r^q-Z_r)\]
is injective. Then we argue, just as in the case above, that the function $g(X)\triangleq \prod_{j=0}^{t-1} \sigma_{i_j}(X)$ expands to a sum of distinct monomials with coefficients ranging over the product set $\hat{A}=\sigma_{i_0}(A)\cdots \sigma_{i_{t-1}}(A)$ where $\sigma(A)=\condset{\sigma(a)}{a\in A}$. The theorem of~\cite{HLX02} is applied $t$ times to conclude that $\hat{A}$ spans $\Fn$ over $\F$. Now we apply \expref{Claim}{clm:nonconstant} and get that the first coordinate of $g(X)$ (viewing $g(X)$ as a tuple of $n$ polynomials over $\F$) is a non-constant function. Details are provided in \expref{Section}{sec:main}.

\paragraph{From dispersers to extractors}
This part is based on the work of Gabizon and Raz~\cite{GR08} and uses
%
an important theorem of Weil~\cite{Weil48}.
This theorem implies the following. Suppose we evaluate a polynomial $g\in \F[Z_1,\ldots,Z_r]$ of small-enough degree $\deg(g)<\sqrt q $
on a uniformly random sample in $\F^r$ and then take the first bit of this evaluation (when viewing it as a vector over $\mathbb{F}_2$).
%
%
%
Then, this bit will either be constant (in which case we then say
%
%
that $g$ is ``degenerate,'' or close to the uniform distribution.
%
 Assuming our source is low-degree and the
%
%
order $q$ of the field
is sufficiently large, we %
can argue that $\deg(E_0(X))<\sqrt q$ because $X$ is
low-degree by assumption and $E_0$ is low-degree by construction. So to apply Weil's Theorem and get an extractor we only need to ensure that we have in hand
a non-degenerate polynomial. Alas, we have relatively little control over the polynomial source so
we %
need to transform it somehow into a non-degenerate one in a
 black-box manner.
Here we apply another observation,
%
%
proved by Swastik Kopparty,
which says that $(E_0(X))^v$ is non-degenerate for
odd\footnote{For characteristic $p>2$ the criteria for $v$ is a bit different:
we need
%
$p\nmid v$.} $v>2$.
 This part is explained in \expref{Section}{sec-weil}. So we take $E_1(Y)$ to be the first\footnote{In fact, we can  output several bits. See \expref{Section}{subsec:weil} for details.}  bit of $Y^3$ and using this observation and Weil's Theorem conclude
 that $E_1(E_0(X))$ is close to uniform. Analysis of the resulting extractor is given in the appendix.

\section{Preliminaries}\label{sec-prelim}
\paragraph{Notation:}
When we discuss identities between polynomials we
only mean identities as \emph{formal polynomials}.
We will frequently alternate between viewing $\x\in \F^n$ as an element of
either $\F^n$ or the field $\Fn$.
When we do this we assume it is using an implicit bijective map $\phi:\F^n\to \Fn$ that is
an isomorphism of vector spaces.
That is, $\phi(t_1\cdot a_1+t_2\cdot a_2)= t_1\cdot \phi(a_1) + t_2\cdot \phi(a_2)$
for any $t_1,t_2\in \F$ and $a_1,a_2\in \F^n$.
Such $\phi$ is efficiently computable using standard representations of $\Fn$.
(For details see for example the book of Lidl and Niederreiter~\cite{LN94}.)
For a set $\Omega$ we denote by $U_{\Omega}$ the uniform distribution
%
over $\Omega$.

\subsection{Weil bounds for additive character sums}\label{subsec:weil}
The seminal work of Weil~\cite{Weil48} on the ``Riemann hypothesis for
curves over finite fields'' implies very useful bounds on character
sums. As we will see in this section, these bounds enable us to
extract randomness from certain ``low-degree distributions.''

 For background on characters of finite fields see~\cite{Schmidt76} or Section~3.2 of~\cite{GR08}.
The following version of the Weil bound was proved by Carlitz and Uchiyama~\cite{CU57}.
\begin{theorem}[Weil-Carlitz-Uchiyama bound]\label{thm:weil_additive}
Let $q=p^{\ell}$ for prime $p$ and an integer $\ell$.
Let $\psi$ be a non-trivial additive character of $\F$ (that is,
not identically $1$). Let $f(Z)$ be a polynomial in $\F[Z]$ of
degree $d$. Suppose that $f$ is not of the form $h^p + h + c$ for any $h\in \F[Z]$ and $c\in \F$.
 Then
\[\left|\sum_{z\in \F} \psi(f(z))\right| \leq (d-1)\cdot q^{1/2}\,.\]
\end{theorem}

We require the following generalization of Vazirani's XOR Lemma from Rao~\cite{Rao07}, appearing there as Lemma 4.2.
\begin{lemma}[Rao's XOR lemma] \label{lem-rao}
Let $X$ be a distribution on a finite abelian group $G$ s.t.
$|\Exp(\psi(X))|\leq \eps$ for any non-trivial character
$\psi$ of $G$. Then $X$ is $\eps \cdot \sqrt {|G|}$-close to
uniform on $G$.
\end{lemma}

The above lemma implies it suffices to bound additive character sums of a distribution over $\F$ in order to extract randomness. This is formalized in \expref{Lemma}{lem-bits_from_small_bias} below. To state the lemma we first define how to extract a few entries of an element in $\mathbb{F}_{p^\ell}$.

\begin{definition}[Prefix projection]\label{dfn-prefix_function}
Let $q=p^{\ell}$ for prime $p$ and an integer $\ell$.
Fix an isomorphism between $\F$ and $\mathbb{F}_p^{\ell}$
and view $x\in \F$ as $(x_1,\ldots,x_{\ell})\in \mathbb{F}_p^{\ell}$.
Fix an integer $m\leq \ell$. We define the \emph{prefix projection function}
$E_m:\F \to \mathbb{F}_p^m$ by $E_m(x) = E_m( (x_1,\ldots,x_{\ell})) \triangleq (x_1,\ldots,x_m)$.
\end{definition}
\begin{lemma}[XOR lemma for prefix projections]\label{lem-bits_from_small_bias}
Let $q=p^{\ell}$ for prime $p$ and an integer $\ell$.
 Let $X$ be a distribution on $\F$ such that
 $|\Exp(\psi(X))|\leq \eps$ for any
 non-trivial additive character $\psi$ of $\F$.
 Then  $E_m(X)$ is $p^{m/2}\cdot \eps$-close to uniform.
\end{lemma}
\begin{proof}
We claim that a function of the form $\psi'(a)\triangleq \psi(E_m(a))$ where
$\psi$ is a character of $\mathbb{F}_p^m$, is a character of $\F$:
Let $\omega \in \mathbb{C}$ be a primitive
%
$p$-th root of unity.
The additive characters of $\F$ are exactly the functions
$\psi:\F\to \mathbb{C}$ of the form $\psi(a)=  \omega^ {T(a)}$
where $T:\F\to \mathbb{F}_p$ is an $\mathbb{F}_p$-linear
function and $T(a)$ is interpreted as an integer in
$\{0,\ldots,p-1\}$. In particular, this includes such functions where $T$ only looks
at the first $m$ coordinates of $a$ (recall that we identify $\F$ with
$\mathbb{F}_p^{\ell}$);
and such functions in turn, are exactly those of the form
$\psi(E_m(a))$ where
$\psi$ is a character of $\mathbb{F}_p^m$.
Hence, from the assumption of the lemma $|\Exp(\psi(E_m(X)))|\leq
\eps$ for any non-trivial additive character of  $\mathbb{F}_p^m$.
{}From \expref{Lemma}{lem-rao}, we have that $E_m(X)$ is $p^{m/2}\cdot
\eps$-close to uniform.

\end{proof}

Summing up the previous results we reach the statement that will be later used in analyzing our extractors.

\begin{corollary}[Weil-Carlitz-Uchiyama for prefix projections]\label{cor-bits_from_weil}
Let $q=p^{\ell}$ for prime $p$ and an integer $\ell$.
Let $f(Z)$ be a polynomial in $\F[Z]$ of
degree $d$. Suppose that $f$ is not of the form $h(Z)^p + h(Z) + c$ for any $h(Z)\in \F[Z]$ and $c\in \F$.
Then $E_m(f(U_{\F}))$ is $p^{m/2}\cdot d/\sqrt q$-close to uniform.
\end{corollary}
\begin{proof}
Follows immediately from \expref{Theorem}{thm:weil_additive}
and \expref{Lemma}{lem-bits_from_small_bias}.
\end{proof}




\subsection{Dimension expansion of products}\label{sec:dim_exp}
Recall that $\Fn$ is a vector space
over $\F$ isomorphic to $\F^n$.
For a set $A\subseteq \mathbb{F}_{q^n}$ we denote by $\dim(A)$
the dimension of the $\F$-span of $A$.
For sets $A,B\subseteq \Fn$ let \[A\cdot B\triangleq \{a\cdot b\mid a\in A, b\in B\}\,.\]
Hou, Leung and Xiang~\cite{HLX02} show that such products expand in dimension.
 The following theorem is a corollary of Theorem 2.4 of~\cite{HLX02}.
\begin{theorem}[Dimension expansion of products]\label{thm-hlx}
Let $\F$ be any field, and let $n$ be prime.\footnote{The theorem
of~\cite{HLX02}
works also for non-prime $n$ in which case the inequality involves
the %
order of a certain subfield of $\Fn$.}
Let $A$ and $B$ be non-empty subsets of $\mathbb{F}_{q^n}$
such that $A,B\neq \{0\}$. Then
\[\dim(A\cdot B)\geq \min\{n,\dim(A)+\dim(B) -1\}\,.\]
In particular,
if $A_1,\ldots, A_m$ are non-empty subsets of $\mathbb{F}_{q^n}$
such that for all $1\leq i \leq m $, $dim(A_i)\geq k$ for some $k\geq 1$. Then
\[\dim(A_1\cdots A_m)\geq  \min\{n,k\cdot m - (m-1)\}\,.\]

\end{theorem}
\begin{remark}
The definition of $A\cdot B$ is somewhat different from that in~\cite{HLX02} where
it is defined only for subspaces, and as the \emph{span} of all possible products.
The definition above will be more convenient for us. It is easy to see
that Theorem 2.4 of~\cite{HLX02} %
%
%
%
is equivalent to
the theorem above with our definition.
Still, we give a self-contained proof.\footnote{Also, see Section 3.2 of~\cite{DG10}
for a self-contained proof using the definition of~\cite{HLX02}.}
\end{remark}
\begin{proof}
First we note that it is enough to prove the theorem for linear subspaces $A$ and
$B$ of dimension at least one:
Given arbitrary sets $A$ and $B$, let $A'\triangleq\spn(A)$ and $B'\triangleq\spn(B)$.
If $A$ and $B$ both contain a non-zero element (as required in the theorem), then $A'$ and $B'$ are linear subspaces
of dimension at least one. So we have that
\[\dim(A'\cdot B')\geq \min\{n,\dim(A')+\dim(B') -1\} = \min\{n,\dim(A)+\dim(B) -1 \}\,.\]
Now, we observe that $\spn(A'\cdot B')\subseteq \spn(A\cdot B)$: An element of $A' \cdot B'$ has the
form
\[(\sum_i t_i\cdot a_i )\cdot( \sum_j s_j\cdot b_j)=\sum_{i,j} t_i\cdot s_j\cdot a_i\cdot b_j\,,\]
where $a_i\in A,b_j\in B$ and $t_i,s_j\in \F$. This is obviously in $\spn(A\cdot B)$.
So $A'\cdot B'\subseteq \spn(A\cdot B)$, and this implies $\spn(A'\cdot B')\subseteq \spn(A\cdot B)$.
Therefore, the equation above implies
\[\dim(A\cdot B)\geq \min\{n,\dim(A)+\dim(B) -1 \}\,.\]
We now turn to proving the theorem for linear subspaces $A$ and $B$ of dimension at least one.
We proceed by induction on $\dim(A)$.  As a base, observe that the
result holds trivially when $\dim(A) = 1$.  For the inductive step, we
may then assume that $\dim(A) >1$.  We may also assume that $B \neq
\Fn$ as the theorem is immediate in this case.

Note that we may freely replace $A$ by $g\cdot A$ (or $B$ by
$g\cdot B$) for any $g \in \Fn$ as this has no effect on $\dim(A)$, $\dim(B)$, or
$\dim(A \cdot B)$.  By this operation, we may assume that $1 \in A \cap
B$.
Since $\dim(A) > 1$, we may choose $a \in A \setminus \F$.  Let $\ell$ be the
smallest nonnegative integer so that $a^{\ell} \not\in B$. Note that such $\ell$ exists
since $\Fn = \spn(1,a,a^2,\ldots,a^{n-1})$ for any $a \in \Fn \setminus \F$
as there are no non-trivial subfields $\F\subsetneq \mathbb{K} \subsetneq \Fn$ when $n$ is prime, and  $B \neq \Fn$. Furthermore, $\ell > 0$ by the assumption that $1 \in B$.  Next, replace $B$
by the set $a^{-(\ell -  1)}\cdot B$.  It now follows that $1 \in B$ and $a
\not\in B$, so $A \cap B$ is a proper nonempty subset of $A$.
In particular, $1\leq \dim(A\cap B) < \dim(A)$.


Consider the $\F$-linear subspaces $A \cap B$ and $A+B$ and observe
that $(A \cap B) \cdot (A + B) \subseteq \spn(A \cdot B)$.
 The next
equation follows from this and the induction hypothesis applied to $A \cap B$ and $A
+B$.
\begin{align*}
 \dim(A\cdot B) &\ge \dim( (A \cap B) \cdot (A + B)) \\
 &\ge \min\{n,\dim(A\cap B)+\dim(A+B) -1\} \\
 &= \min\{n,\dim(A) + \dim(B) - 1\}\,.
\end{align*}
This completes the proof.
\end{proof}
\subsection{Frobenius automorphisms of \texorpdfstring{$\F$}{F}}\label{subsec:frobenius}
Let $q=p^{\ell}$ for prime $p$ and let $i\geq 0$ be an integer.
Raising to power $p^i$ in $\F$ is known as a Frobenius automorphism of $\F$ over $\mathbb{F}_p$ and will play an important role. We record two useful and well-known properties of this automorphism that will be used in our proofs.
\begin{itemize}
\item{\bf Linearity:} $\forall a,b \in \F$, $(a+b)^{p^i} = a^{p^i} + b^{p^i}$.
\item{\bf Bijection:} The map $x\to x^{p^i}$ over $\F$ is bijective. In particular,
for $c\in \F$, $c^{1/p^i}$ is always (uniquely) defined.
\end{itemize}

A useful fact following from these properties is that ``taking the
%
$p$-th
power'' of a set does not change its dimension.
\begin{claim}[Dimension preservation]\label{claim-frobenius_has_dim}
Let $q=p^{\ell}$ from prime $p$ and an integer $\ell$.
For an integer $i\geq 1$ and a set $A\subseteq\Fn$
let $A^{p^i}\triangleq \{a^{p^i}\mid a\in A\}$.
Then $\dim(A)=\dim(A^{p^i})$.
\end{claim}
\begin{proof}
Let $\{a_1,\ldots, a_k\}\subseteq A$ be a basis for the $\F$-span of $A$.
Choose any $c_1,\ldots,c_k\in \F$ that are not all zero.
 Then,
\[\sum_{j=1}^k c_j \cdot a_j^{p^i} = \left(\sum_{j=1}^k c_j^{1/p^i} \cdot a_j\right)^{p^i}\neq 0\,.\]
Thus $\{a_1^{p^i},\ldots, a_k^{p^i}\}$ are independent over $\F$ and therefore $\dim(A^{p^i})\geq \dim(A)$.
The reverse inequality is similar.
\end{proof}


 \section{The main construction}\label{sec:main}
As before, we use $r$ to denote the number of inputs of $f(Z_1,\ldots,Z_r)\in\M{n}{k}{d}$.
We denote by $\D$ the product set $\{0,\ldots,d\}^r$.
We use bold letters to denote vectors in $\F^r$.
For example, $\zi = (Z_1,\ldots,Z_r)$.
For an element $S=(s_1,\ldots,s_r)\in \D$ we use the notation
\[\zi^{S}\triangleq Z_1^{s_1}\cdots Z_r^{s_r}\,.\]
Fix $f=(f_1(\zi),\ldots,f_n(\zi))\in \M{n}{k}{d}$.
For $1\leq j\leq n$, we write
\[f_j(\zi) = \sum_{S \in \D} a_{j,S} \cdot \zi^S\,.\]
With the notation above, for $S\in \D$ let $a_S\triangleq
(a_{1,S},\ldots,a_{n,S})\in \F^n$. Using the isomorphism of the
vectors spaces $\F^n$ and $\Fn$, we can view $a_S$ as an element
of $\Fn$ and write
\begin{equation}
  \label{eq:univariate}
f(\zi) = \sum_{S \in \D} a_S \cdot \zi^S.
\end{equation}
That is, we view $f$ as a multivariate polynomial with coefficients in $\Fn$.
A crucial observation is that when $f$ has large range the coefficients
of $f$ have large dimension.
\begin{lemma}[Large range implies large span]\label{lem-spanofcoeff}
Let $f\in \M{n}{k}{d}$. As in (\ref{eq:univariate}), write $f(\zi)=\sum_{S\in \D} a_S\cdot
\zi^S$ where $a_S\in \Fn$. Then $\dim \{a_S\}_{ S\in \D\setminus \zero}\geq  k$.
\end{lemma}
\begin{proof}
The range of $f$ over inputs in $\F^r$ is contained in an affine
shift  of the $\F$-linear span of
$\{a_S\}_{S\in \D\setminus \zero}$. Since this range is of size at
least $q^k$, we must have $\dim \{a_S\}_{S\in \D\setminus
\zero}\geq  k$.
\end{proof}

%
%
%
%
%

A simple but crucial observation from~\cite{DG10} is that a polynomial with coefficients in $\Fn$ whose
non-constant coefficients span $\Fn$ over $\F$ can be ``projected'' to a non-constant polynomial with coefficients
in $\F$. We formalize this in the definition and lemma below.
%
%
%
%
\begin{definition}[Full-span polynomial]
  We say that a polynomial $G\in \Fn[\zi]=\Fn[Z_1,\ldots,Z_r]$ has \emph{full span}
  if the coefficients of the non-constant monomials of $G$
  span $\Fn$ over $\F$.

\end{definition}
\begin{lemma}[Disperser for full-span polynomials]\label{lem-can_project_expanding}
Suppose $G\in \Fn[\zi]$ has full span. Let $T:\Fn \to \F$ be a non-trivial $\F$-linear mapping.
Then $T(G(\zi))$, as a function from $\F^r$ to $\F$, agrees with a non-constant polynomial in $\F[\zi]$
whose total and individual degrees are at most those of $G$.
\end{lemma}
\begin{proof}
We write  $G(\zi)= \sum_{S\in {\cal R}} a_S\cdot
\zi^S$ for  $a_S\in \Fn$, where ${\cal R}\subset {\mathbb N}^r$ denotes the set of tuples corresponding
to the monomials of $G$.
For every $\mathbf{x}=(x_1,\ldots,x_r) \in \F^r$, we have
\[T(G(\mathbf{x})) = T\left( \sum_{S\in {\cal R}} a_S\cdot
\mathbf{x}^S\right) = \sum_{S\in {\cal R}}T(a_S)\cdot \mathbf{x}^S\,,\]
where the last inequality used the $\F$-linearity of $T$.
Thus $T(G(\zi))$ agrees on all inputs in $\F^r$ with the polynomial
$F(\zi)\triangleq \sum_{S\in {\cal R}}T(a_S)\cdot \mathbf{\zi} ^S$ which
is in $\F[\zi]$.
The full span of $G$ means that
  $\dim \{a_S\}_{S\in {\cal R}\setminus
\zero}=n$. Since $T$ is a nontrivial linear map there is some $S \in {\cal R}$ such that $T(a_S)\neq 0$ and $S\neq \mathbf{0}$
and so $F$ is a non-constant polynomial.
As the monomials with non-zero coefficients in
$F$ are a subset of the monomials with non-zero coefficients in $G$,
it is clear that
%
%
the total and individual degrees of $F$
%
are at most those of $G$.
\end{proof}

%
%
%
%
%
%
The previous lemma implies that to construct a disperser for
polynomial sources it suffices to produce a function that increases
the span of low-degree polynomials.  We do this in the next
theorem which is of paramount importance to this paper.
%
\begin{theorem}[Product of distinct Frobenius automorphisms increases span]\label{thm-spanning_poly}
Fix a prime power $q=p^{\ell}$. Fix integers $k\leq n$ and $d<s$ such that $n$ is prime
and $s$ is a power of $p$. (In particular, raising to power $s^i$ is a Frobenius automorphism of $\F$ over $\mathbb{F}_p$.)
Let $u= \lceil (n-k)/(k-1)\rceil $.
Then for any $f(Z_1,\ldots,Z_r)\in \M{n}{k}{d}$, the polynomial
\[f^{1+s+s^2+\cdots +  s^u}(Z_1,\ldots,Z_r)=f(Z_1,\ldots,Z_r)\cdot f^s(Z_1,\ldots,Z_r)\cdots f^{s^u}(Z_1,\ldots,Z_r)\] has full span.
%
%
%
\end{theorem}
\begin{proof}
Fix $f\in \M{n}{k}{d}$.
As in (\ref{eq:univariate}), write $f(\zi)=\sum_{S\in \D} a_S
\cdot\zi^S$ with $a_S\in \Fn$.
\begin{align*}
f^{1+s+s^2+\cdots +s^u}(\zi) &=
\left(\sum_{S\in \D} a_S \cdot \zi^S\right)^{1+s+s^2+\cdots +s^u}=
\prod_{i=0}^u \left(\sum_{S\in \D} a_S \cdot \zi^S\right)^{s^i}\,.\\
\intertext{In what follows we use the notation $S_i=(S_{i,1},\ldots,S_{i,r})$ and $S_i\cdot s^i=(S_{i,1}\cdot s^i,\ldots,S_{i,r}\cdot s^i)$. Using the linearity of Frobenius automorphisms we continue the derivation and get}
%
%
&= \prod_{i=0}^u \left(\sum_{S\in \D} a_S^{s^i} \cdot \zi^{S\cdot s^i}\right)
=\sum_{S_0,\ldots,S_u \in \D} \prod_{i=0}^u a_{S_i}^{s^i}\cdot\prod_{i=0}^u \zi^{S_i\cdot s^i}\\
&=\sum_{S_0,\ldots,S_u \in \D} \prod_{i=0}^u a_{S_i}^{s^i}\cdot
 \prod_{i=0}^u \prod_{j=1}^r Z_j^{S_{i,j}\cdot s^i} =
 \sum_{S_0,\ldots,S_u \in \D} A_{S_0,\ldots,S_u} \cdot M_{S_0,\ldots,S_u}(\zi)\,,\end{align*}
where 
\[
A_{S_0,\ldots,S_u}=\prod_{i=0}^u a_{S_i}^{s^i}\quad\text{and}\quad
M_{S_0,\ldots,S_u}(\zi)=\prod_{i=0}^u \prod_{j=1}^r Z_j^{S_{i,j}\cdot
  s^i}\,.
\]
The crucial observation is that if $(S_0,\ldots, S_u)$ and $(S'_0,\ldots, S'_u)$ are two distinct tuples of elements of $\D$ then the monomials $M_{S_0,\ldots, S_u}(\zi)$
and $M_{S'_0,\ldots, S'_u}(\zi)$ are distinct as well:
Consider $j\in\{1,\ldots,r\}$ such that $S_{i,j}\neq S'_{i,j}$ for some $0\leq i \leq u$.
Then $Z_j$ is raised to power $\sum_{i=0}^u S_{i,j}\cdot s^i$ in $M_{S_0,\ldots, S_u}(\zi)$
and to power $\sum_{i=0}^u S'_{i,j}\cdot s^i$
in $M_{S'_0,\ldots, S'_u}(\zi)$.
These powers are different as for all $0\leq i \leq u$,
$S_{i,j},S'_{i,j}\leq d<s$, and there is only one way to write an integer in base $s$ with ``coefficients'' smaller than $s$.
Define    %
\[
A\triangleq \{A_{S_0,\ldots,S_u}\mid S_0,\ldots,S_u\in
\D\setminus\zero\}\,.
\]
 For $0\leq i \leq u$,
define    %
\[
B^{s^i}\triangleq \{a_S^{s^i}\mid  S\in \D\setminus \zero\}\,.
\]
Note that $A = B^{s^0}\cdots B^{s^u}$.
For all $0\leq i \leq u$, by \expref{Lemma}{lem-spanofcoeff} and \expref{Claim}{claim-frobenius_has_dim}
we have $\dim(B^{s^i})\geq k$.
Therefore, by \expref{Theorem}{thm-hlx} we get
\[\dim(A)\geq \min\{n, k\cdot (u+1) - u\}=n\,.\]
 Our theorem follows by noticing that the coefficients of the non-constant monomials in  $f^{1+s+s^2+\cdots +  s^u}$
 contain the set $A$, hence $f^{1+s+\cdots+s^u}$ has full span.
\end{proof}

Combining the lemma and theorem above we ``project'' into $\F$ and get a non-constant polynomial with
coefficients in $\F$.
\begin{theorem}\label{thm-main_const}
Fix a prime power $q=p^{\ell}$. Fix integers $k\leq n$ and $d<s$ such that $n$ is prime
and $s$ is a power of $p$. Fix a non-trivial $\F$-linear map $T:\Fn\to \F$.
Let $u= \lceil (n-k)/(k-1)\rceil $.
Define $P:\Fn\to \F$ by $P(x) \triangleq T(x^{1+s+s^2+\cdots +  s^u})$.
Fix any $f(Z_1,\ldots,Z_r)\in \M{n}{k}{d}$ of total degree $D$. Then $P(f(\zi))$, as a function
on $\F^r$, agrees with a non-constant polynomial in $\F[\zi]$ of total degree
at most $D\cdot (1+ s+ s^2 +\cdots + s^u)< D\cdot s^{u+1}$
and individual degree at most $d\cdot(1+ s+ s^2 +\cdots + s^u) = d\cdot (s^{u+1}-1)/(s-1)$.
\end{theorem}
\begin{proof}
Follows immediately from \expref{Lemma}{lem-can_project_expanding} and \expref{Theorem}{thm-spanning_poly}.
\end{proof}

An immediate corollary is a construction of a ``non-Boolean disperser'' for polynomial sources.
\begin{corollary}\label{cor-disp}
Fix a prime power $q=p^{\ell}$. Fix integers $k\leq n$ and $d<s$ such that $n$ is prime
and $s$ is a power of $p$. Fix a non-trivial $\F$-linear map $T:\Fn\to \F$.
Let $u= \lceil (n-k)/(k-1)\rceil $.
Define $P:\Fn\to \F$ by $P(x) \triangleq T(x^{1+s+s^2+\cdots +  s^u})$.
Assume that $q> d\cdot (s^{u+1}-1)/(s-1)$.
Then, for any $f(Z_1,\ldots,Z_r)\in \M{n}{k}{d}$ we have that $P(f(\zi))$ is  a non-constant function
from $\F^r$ into $\F$.
\end{corollary}
\begin{proof}
Follows immediately from \expref{Theorem}{thm-main_const} by noticing that if
$P(f)$ agrees with a non-constant polynomial whose individual degrees are smaller than $q$,
then it is a non-constant function from $\F^r$ into $\F$.
\end{proof}

\section{A useful criteria for the Weil bound}\label{sec-weil}
To get our main result we shall apply the Weil-Carlitz-Uchiyama bound for prefix projections (\expref{Corollary}{cor-bits_from_weil}) to
a certain polynomial $f\in \F[Z]$, and so we have to ensure that $f$ is not of the
``degenerate'' form $h^p+h+c$ precluded by that bound. The common way
to do this is to require $\gcd(\deg(f),p) =1$ (cf.,~\cite{GR08,DG10}). However we have less control
%
over
the degree of the polynomial $f$ we need to work with.
For this reason, the following lemma will be very helpful to us.
It gives us a simple way to ``alter'' $f$ and get a polynomial that is not of the
form $h^p+ h+c$.
The proof of the following lemma was shown to us by Swastik Kopparty.
\begin{lemma}[Criteria for non-degenerateness]\label{lem-powerhelps}
Let $q=p^{\ell}$ for prime $p$ and let $v\geq 2$ be an integer such that $p\nmid v$.
Let $f\in \F[Z]$ be a non-constant polynomial. If $f$ is of the form $g^v$ for some $g\in \F[Z]$, it
is not of the form $h^p +h + c$ for any $h\in \F[Z]$ and $c\in \F$.
\end{lemma}
\begin{proof}
Suppose by way of contradiction there exists $f\in\F[Z]$ of degree $d\geq 1$
such that
\[f = g^v= h^p + h + c\]
 for some $g,h\in \F[Z]$ and $c\in \F$.
 Fix such an $f$ with minimal degree $d\geq 1$.
It follows that $\deg(g)= d/v$ and $\deg(h) = d/p$.
Taking a derivative in $\F[Z]$ of all 3 parts of the above equation we get
\[f'(Z) =  v\cdot g^{v-1}(Z)\cdot g'(Z)= h'(Z)\,,\]
where in the rightmost part we used the fact that the derivative of $h^p$ is zero.
Notice that $v\neq 0$ in $\F$ since $p\nmid v$.
%
%
If $g'\not\equiv 0$ then this implies $\deg(h') \geq (v-1)\cdot
\deg(g)= d \cdot{(v-1)}/{v}$.
%
%
%
%
%
%
%
%
 But $\deg(h') < d/p \le d\cdot {(v-1)}/{v}$.
 (For the last inequality we use $p\geq 2$ and $v\geq 2$.)
%
So $g'$ and $h'$ are the zero polynomial.
It is not hard to see that this implies that all powers in
$g$ and $h$ are multiples of $p$.
So $g= g_1^p$ and $h=h_1^p$ for some $g_1,h_1\in \F[Z]$.
We now have
\[f = (g_1^p)^v = (h_1^p)^p  + h_1^p + c\,.\]
This implies
\[ g_1^v = h_1^p + h_1 + c^{1/p}\,.\]
%
%
 (Recall that a $p$-th root always exists in $\F$.)
Since $g_1^v$ has positive degree smaller than $\deg(f)=d$, this contradicts the minimality of $d$ and proves the lemma.
\end{proof}
\noindent

Reducing the multivariate case to the univariate case, we get the version of the Weil bound we need.
\begin{lemma}\label{lem-multivariate_weil}
Let $q=p^{\ell}$ for a prime $p$ and integer $\ell>0$. Let $f(Z_1,\ldots,Z_r)\in \F[Z_1,\ldots,Z_r]$ be a
non-constant polynomial of total degree $d<q$.
Assume that $f= g^v$ for an integer $v\geq 2$ with
$p\nmid v$ and some $g\in \F[Z_1,\ldots,Z_r]$. Let $m<\ell $ be a positive integer.
Then $E_m(f(U_{\F^r}))$ is $\eps$-close to uniform for $\eps=p^{m/2}\cdot d\cdot q^{-1/2}$.
\end{lemma}
\begin{proof}
We note first that there must be an $a= (a_1,\ldots,a_r)\in \F^r$ such that the univariate ``line restriction'' polynomial
\[f_a(Z)\triangleq f(a\cdot Z ) = f(a_1\cdot Z,\ldots,a_r\cdot Z)\] has degree \emph{exactly} $d$: The coefficient
of $Z^d$ in $f_a$ is $f^d(a)$ where $f^d$ is the $d$-homogeneous part of $f$, \ie, the sum of monomials of degree
exactly $d$ in $f$. By the Schwartz-Zippel lemma as $d<q$, there is an $a\in \F^r$ such that $f^d(a)\neq 0$
and therefore $f_a(Z)$ has degree $d$.
Fix such an $a\in \F^r$.
It follows that for all $b = (b_1,\ldots,b_r) \in \F^r$,
\[f_{a,b}(Z)\triangleq
f(a\cdot Z+ b) = f(a_1\cdot Z + b_1,\ldots, a_r\cdot Z + b_r)\]
%
%
is non-constant, as the coefficient
%
of $Z^d$ in $f_{a,b}$ is also $f^d(a)$.
 Furthermore, for any $b\in \F^r$
\[f_{a,b}= f(a_1\cdot Z + b_1,\ldots, a_r\cdot Z + b_r)=
g^v(a_1\cdot Z + b_1,\ldots, a_r\cdot Z + b_r)\,,\]
and so $f_{a,b}$ is a
%
$v$-th
power of a polynomial in $\F[Z]$, and so
by \expref{Lemma}{lem-powerhelps} is not of the form $h^p + h + c$ for any $h\in \F[Z]$ and $c\in \F$.
As the distribution $f(U_{\F^r})$ is a convex combination of the distributions $f_{a,b}(U_{\F})$
for the different ``shifts'' $b\in\F^r,$ the claim now follows from the Weil-Carlitz-Uchiyama bound for prefix projections (\expref{Corollary}{cor-bits_from_weil}).
\end{proof}


%
%
\section{A polynomial-source extractor}\label{sec:extractor}

We can now state and prove our main technical theorem, which immediately implies our main theorem on extractors for polynomial sources (\expref{Theorem}{thm-ext_main_intro}).

\begin{theorem}[Main--Extractors, parameterized version]\label{thm-ext_main}
Fix a field $\F$ of characteristic $p$, integers $d,D,2\leq k\leq n$ where $n\geq 25$,
and a positive integer $m<1/2\cdot \log_p q$.
Let $\alpha=3 D \cdot (p\cdot d)^{\frac{1.2\cdot n-k}{k-1} +2} $. Assume that   $q\geq 2\cdot  \alpha^2$.
There is an explicit \polyext{k}{d}{D}{\eps}  $E:\F^n\to \mathbb{F}_p^m$ with error
$\eps =  p^{m/2}\cdot \alpha \cdot q^{-1/2}$.
\end{theorem}

\expref{Theorem}{thm-ext_main_intro} follows
from the previous theorem by noticing that for $4\leq k \leq n$,
\[\frac{1.2\cdot n-k}{k-1} +2 \leq 3n/k\,.\]

\begin{proof}[Proof of \expref{Theorem}{thm-ext_main}]
Choose a prime $n\leq n'\leq 1.2\cdot n$ (which always exists for $n\geq 25 $ according to Nagura's improvement of the
Bertrand-Chebychev Theorem~\cite{Na52}).
Given $f(Z_1,\ldots,Z_r)\in \M{n}{k}{d}$ of total degree $D$ we think of $f$ as an element of $ \M{n'}{k}{d}$
by padding its output with zeros.
Let $s$ be the smallest power of $p$ greater than $d$.
Note that $s\leq p\cdot d$.
Let $P:\F^{n'}\to \F$ be the polynomial in \expref{Theorem}{thm-main_const} using $s$ as above.
If $p=2$ let $v=3$ and otherwise let $v=2$.
Let $E:\F^n\to \mathbb{F}_p^m$ be defined as $E(\x)\triangleq E_m(P^v(\x))$.
{}From \expref{Theorem}{thm-main_const} we conclude that $P(f(\zi))$ is non-constant of degree at most $D\cdot s^{u+1}$
where 
\[
u= \lceil (n'-k)/(k-1)\rceil \leq  \frac{ 1.2\cdot n-k}{k-1} + 1\,.
\]
Hence, from \expref{Lemma}{lem-multivariate_weil} we see that $E_m(P^v(f(U_{\F^r})))$ is $\eps$-close to uniform
for
\[\eps = p^{m/2}\cdot v\cdot D\cdot s^{u+1}\cdot q^{-1/2}\leq
p^{m/2}\cdot 3 D \cdot (p\cdot d)^{\frac{1.2\cdot n-k}{k-1} +2}\cdot q^{-1/2}
=p^{m/2}\cdot \alpha\cdot q^{-1/2}\,.\qedhere\]
\end{proof}

%
%
%
%
%
\providecommand{\bibhead}[1]{}
%
\expandafter\ifx\csname pdfbookmark\endcsname\relax%
  \providecommand{\tocrefpdfbookmark}{}
\else\providecommand{\tocrefpdfbookmark}{%
   \hypertarget{tocreferences}{}%
   \pdfbookmark[1]{References}{tocreferences}}%
\fi

\tocrefpdfbookmark
\begin{thebibliography}{10}

\bibitem{BG12}\bibhead{BG12}
{\sc Eli {Ben-Sasson} and Ariel Gabizon}: Extractors for polynomials sources
  over constant-size fields of small characteristic.
\newblock In {\em Proc. 16th Internat. Workshop on Randomization and
  Computation (RANDOM'12)}, pp. 399--410. Springer, 2012.
\newblock [\epfmtdoi{10.1007/978-3-642-32512-0\_34}]

\bibitem{BHRVW}\bibhead{BHRVW}
{\sc Eli {Ben-Sasson}, Shlomo Hoory, Eyal Rozenman, Salil Vadhan, and Avi
  Wigderson}: Extractors for affine sources.
\newblock Unpublished Manuscript, 2001.

\bibitem{BK09}\bibhead{BK09}
{\sc Eli {Ben-Sasson} and Swastik Kopparty}: Affine dispersers from subspace
  polynomials.
\newblock {\em SIAM J. Comput.}, 41(4):880--914, 2012.
\newblock Preliminary version in
  \href{http://dx.doi.org/10.1145/1536414.1536426}{STOC'09}.
\newblock [\epfmtdoi{10.1137/110826254}]

\bibitem{blum}\bibhead{blum}
{\sc Norbert Blum}: A {B}oolean function requiring $3n$ network size.
\newblock {\em Theoret. Comput. Sci.}, 28(3):337--345, 1983.
\newblock [\epfmtdoi{10.1016/0304-3975(83)90029-4}]

\bibitem{Bour05}\bibhead{Bour05}
{\sc Jean Bourgain}: On the construction of affine extractors.
\newblock {\em Geometric and Functional Analysis}, 17(1):33--57, 2007.
\newblock [\epfmtdoi{10.1007/s00039-007-0593-z}]

\bibitem{CU57}\bibhead{CU57}
{\sc Leonard Carlitz and Sabur{\^o} Uchiyama}: Bounds for exponential sums.
\newblock {\em Duke Math. J.}, 24(1):37--41, 1957.
\newblock [\epfmtdoi{10.1215/S0012-7094-57-02406-7}]

\bibitem{CG88}\bibhead{CG88}
{\sc Benny Chor and Oded Goldreich}: Unbiased bits from sources of weak
  randomness and probabilistic communication complexity.
\newblock {\em SIAM J. Comput.}, 17(2):230--261, 1988.
\newblock Preliminary version in
  \href{http://dx.doi.org/10.1109/SFCS.1985.62}{FOCS'85}.
\newblock [\epfmtdoi{10.1137/0217015}]

\bibitem{DeW11}\bibhead{DeW11}
{\sc Anindya De and Thomas Watson}: Extractors and lower bounds for locally
  samplable sources.
\newblock {\em ACM Trans. Computation Theory}, 4(1):3, 2012.
\newblock Preliminary version in
  \href{http://dx.doi.org/10.1007/978-3-642-22935-0_41}{RANDOM'11}.
\newblock [\epfmtdoi{10.1145/2141938.2141941}]

\bibitem{DK11}\bibhead{DK11}
{\sc Evgeny Demenkov and Alexander~S. Kulikov}: An elementary proof of a $3n -
  o(n)$ lower bound on the circuit complexity of affine dispersers.
\newblock In {\em 36th Internat. Symp. on Mathematical Foundations of Computer
  Science (MFCS'11)}, pp. 256--265, 2011.
\newblock [\epfmtdoi{10.1007/978-3-642-22993-0\_25}]

\bibitem{DG10}\bibhead{DG10}
{\sc Matt DeVos and Ariel Gabizon}: Simple affine extractors using dimension
  expansion.
\newblock In {\em Proc. 25th IEEE Conf. on Computational Complexity (CCC'10)},
  pp. 50--57. IEEE Comp. Soc. Press, 2010.
\newblock [\epfmtdoi{10.1109/CCC.2010.14}]

\bibitem{Dvir09}\bibhead{Dvir09}
{\sc Zeev Dvir}: Extractors for varieties.
\newblock {\em Comput. Complexity}, 21(4):515--572, 2012.
\newblock Preliminary version in
  \href{http://dx.doi.org/10.1109/CCC.2009.7}{CCC'09}.
\newblock [\epfmtdoi{10.1007/s00037-011-0023-3}]

\bibitem{DGW09}\bibhead{DGW09}
{\sc Zeev Dvir, Ariel Gabizon, and Avi Wigderson}: Extractors and rank
  extractors for polynomial sources.
\newblock {\em Comput. Complexity}, 18(1):1--58, 2009.
\newblock Preliminary version in
  \href{http://dx.doi.org/10.1109/FOCS.2007.9}{FOCS'07}.
\newblock [\epfmtdoi{10.1007/s00037-009-0258-4}]

\bibitem{DL11}\bibhead{DL11}
{\sc Zeev Dvir and Shachar Lovett}: Subspace evasive sets.
\newblock In {\em Proc. 44th STOC}, pp. 351--358. ACM Press, 2012.
\newblock See also at \href{http://eccc.hpi-web.de/report/2011/139/}{ECCC}.
\newblock [\epfmtdoi{10.1145/2213977.2214010}]

\bibitem{GR08}\bibhead{GR08}
{\sc Ariel Gabizon and Ran Raz}: Deterministic extractors for affine sources
  over large fields.
\newblock {\em Combinatorica}, 28(4):415--440, 2008.
\newblock Preliminary version in
  \href{http://dx.doi.org/10.1109/SFCS.2005.31}{FOCS'05}.
\newblock [\epfmtdoi{10.1007/s00493-008-2259-3}]

\bibitem{G11}\bibhead{G11}
{\sc Venkatesan Guruswami}: Linear-algebraic list decoding of folded
  {R}eed-{S}olomon codes.
\newblock In {\em Proc. 26th IEEE Conf. on Computational Complexity (CCC'11)},
  pp. 77--85. IEEE Comp. Soc. Press, 2011.
\newblock [\epfmtdoi{10.1109/CCC.2011.22}]

\bibitem{HLX02}\bibhead{HLX02}
{\sc {Xiang-Dong} Hou, Ka~Hin Leung, and Qing Xiang}: A generalization of an
  addition theorem of {K}neser.
\newblock {\em J. Number Theory}, 97(1):1--9, 2002.
\newblock [\epfmtdoi{10.1006/jnth.2002.2793}]

\bibitem{Li11}\bibhead{Li11}
{\sc Xin Li}: A new approach to affine extractors and dispersers.
\newblock In {\em Proc. 26th IEEE Conf. on Computational Complexity (CCC'11)},
  pp. 137--147. IEEE Comp. Soc. Press, 2011.
\newblock [\epfmtdoi{10.1109/CCC.2011.27}]

\bibitem{LN94}\bibhead{LN94}
{\sc Rudolf Lidl and Harald Niederreiter}: {\em Introduction to Finite Fields
  and Their Applications}.
\newblock Cambridge Univ. Press, Cambridge, 1994.

\bibitem{Na52}\bibhead{Na52}
{\sc Jitsuro Nagura}: On the interval containing at least one prime number.
\newblock {\em Proc. Japan Acad.}, 28(4):177--181, 1952.
\newblock [\epfmtdoi{10.3792/pja/1195570997}]

\bibitem{Rao07}\bibhead{Rao07}
{\sc Anup Rao}: An exposition of {B}ourgain's 2-source extractor.
\newblock {\em Electron. Colloq. on Comput. Complexity (ECCC)}, 14(034), 2007.
\newblock \href{http://eccc.hpi-web.de/report/2007/034/}{ECCC}.

\bibitem{Schmidt76}\bibhead{Schmidt76}
{\sc Wolfgang~M. Schmidt}: {\em Equations over Finite Fields: An Elementary
  Approach}.
\newblock Volume 536 of {\em Lecture Notes in Mathematics}.
\newblock Springer, 1976.
\newblock [\epfmtdoi{10.1007/BFb0080437}]

\bibitem{Sha11}\bibhead{Sha11}
{\sc Ronen Shaltiel}: Dispersers for affine sources with sub-polynomial
  entropy.
\newblock In {\em Proc. 52nd FOCS}, pp. 247--256. IEEE Comp. Soc. Press, 2011.
\newblock Full version available at
  \href{http://cs.haifa.ac.il/~ronen/online_papers/online_papers.html}{author's
  home page}.
\newblock [\epfmtdoi{10.1109/FOCS.2011.37}]

\bibitem{viola}\bibhead{viola}
{\sc Emanuele Viola}: Extractors for circuit sources.
\newblock In {\em Proc. 52nd FOCS}, pp. 220--229. IEEE Comp. Soc. Press, 2011.
\newblock See also at \href{http://eccc.hpi-web.de/report/2011/056/}{ECCC}.
\newblock [\epfmtdoi{10.1109/FOCS.2011.20}]

\bibitem{vN}\bibhead{vN}
{\sc John von Neumann}: Various techniques used in connection with random
  digits.
\newblock {\em Applied Math Series}, 12:36--38, 1951.

\bibitem{Weil48}\bibhead{Weil48}
{\sc Andr{\'e} Weil}: On some exponential sums.
\newblock {\em Proc. Nat. Acad. Sci. USA}, 34(5):204--207, 1948.
\newblock \href{https://www.ncbi.nlm.nih.gov/pmc/articles/PMC1079093/}{PNAS}.

\bibitem{Yehud09}\bibhead{Yehud09}
{\sc Amir Yehudayoff}: Affine extractors over prime fields.
\newblock {\em Combinatorica}, 31(2):245--256, 2011.
\newblock [\epfmtdoi{10.1007/s00493-011-2604-9}]

\bibitem{BZ11}\bibhead{BZ11}
{\sc Noga Zewi and Eli {Ben-Sasson}}: From affine to two-source extractors via
  approximate duality.
\newblock In {\em Proc. 43rd STOC}, pp. 177--186. ACM Press, 2011.
\newblock [\epfmtdoi{10.1145/1993636.1993661}]

\end{thebibliography}

